\documentclass[aspectratio=169]{beamer}
\usepackage[utf8]{inputenc}
\usepackage[T2A]{fontenc}
\usepackage[russian]{babel}
\usepackage{amsmath}
\usepackage{amsfonts}
\usepackage{amssymb}
\usepackage{graphicx}
\usepackage{tikz}
\usepackage{hyperref}
\usepackage{subcaption}
\usepackage{color}
\usepackage{booktabs}

% Настройка темы
\usetheme{Madrid}
\usecolortheme{default}

% Настройка цветов
\definecolor{miptred}{RGB}{178,34,52}
\definecolor{miptblue}{RGB}{0,102,153}
\definecolor{darkgreen}{RGB}{0,100,0}

\setbeamercolor{structure}{fg=miptblue}
\setbeamercolor{frametitle}{fg=miptblue}

% Информация о презентации
\title[Инвариантность к сдвигам в CNN]{Методы повышения инвариантности к сдвигам в сверточных нейронных сетях для задач компьютерного зрения}
\subtitle{Выпускная квалификационная работа}
\author{Полев Алексей}
\institute[МФТИ]{Московский физико-технический институт}
\date{2025}

\begin{document}

% Титульный слайд
\begin{frame}
    \titlepage
\end{frame}

% Содержание
\begin{frame}{Содержание}
    \tableofcontents
\end{frame}

\section{Введение}

\begin{frame}{Сверточные нейронные сети (CNN)}
    \begin{block}{Что такое CNN?}
        Класс глубоких нейронных сетей, специально разработанных в первую очередь для обработки изображений
    \end{block}
    
    \begin{columns}
        \begin{column}{0.5\textwidth}
            \textbf{Ключевые компоненты:}
            \begin{itemize}
                \item \textbf{Свертка} — извлечение локальных признаков
                \item \textbf{Пулинг} — уменьшение размерности
                \item \textbf{Активация} — нелинейные преобразования
                \item \textbf{Полносвязные слои} — классификация
            \end{itemize}
        \end{column}
        \begin{column}{0.5\textwidth}
            \textbf{Применения:}
            \begin{itemize}
                \item Классификация изображений
                \item Детекция объектов
                \item Сегментация
                \item Медицинская диагностика
                \item Автономное вождение
                \item Системы безопасности
            \end{itemize}
        \end{column}
    \end{columns}
    

\end{frame}

\begin{frame}{Что такое алиасинг?}
    \begin{block}{Определение}
        \textbf{Алиасинг} — искажение сигнала при нарушении теоремы Найквиста–Шеннона
    \end{block}

    \begin{center}
        \textbf{Теорема Найквиста–Шеннона}
        \[
            f_s \;\ge\; 2\,f_{\text{max}}
        \]
        
        \vspace{0.5cm}
        
        \begin{block}{Применение к CNN}
            В CNN операции с \texttt{stride}$>1$ «отбрасывают» выборки → высокие частоты маскируются под низкие частоты.
            
            \vspace{0.3cm}
            
            \textbf{Результат:} При малом сдвиге входного изображения на 1 пиксель, выходные активации могут кардинально измениться.
        \end{block}
    \end{center}
\end{frame}

\begin{frame}{Иллюстративный пример алиасинга в MaxPool}
    \begin{center}
        \textbf{Иллюстративный пример алиасинга с MaxPool}
        
        \vspace{0.5cm}
        \begin{tikzpicture}[scale=1.0]
            % Оси (с правильными отступами)
            \draw[->] (0,0) -- (9,0) node[right] {\small Пространственная позиция};
            \draw[->] (0,0) -- (0,2.0) node[above] {\small Сигнал};
            
            % Входной сигнал (серая линия с точками) - координаты по вашему описанию
            \draw[thick,gray] plot coordinates {(0.5,0) (1.5,1) (2.5,1) (3.5,0) (4.5,0) (5.5,1) (6.5,1) (7.5,0)};
            
            % Точки входного сигнала на позициях 0.5, 1.5, 2.5, ...
            \foreach \x/\y in {0.5/0, 1.5/1, 2.5/1, 3.5/0, 4.5/0, 5.5/1, 6.5/1, 7.5/0}{
                \fill[gray] (\x,\y) circle (3pt);
            }
            
            % MaxPool без сдвига (синие квадраты) - над четными позициями (0,2,4,6)
            % Берем максимум из окон [0.5,1.5], [2.5,3.5], [4.5,5.5], [6.5,7.5]
            \foreach \x/\y in {0/0, 2/1, 4/0, 6/1, 8/0}{
                \draw[thick,blue,fill=blue!20] (\x-0.12,\y-0.12) rectangle (\x+0.12,\y+0.12);
            }
            
            % MaxPool со сдвигом на 1 (красные ромбы) - над нечетными позициями (1,3,5,7)
            % Берем максимум из окон [1.5,2.5], [3.5,4.5], [5.5,6.5], [7.5,8.5]
            \foreach \x/\y in {1/1, 3/1, 5/1, 7/1}{
                \draw[thick,red,fill=red!20] (\x-0.15,\y) -- (\x,\y+0.15) -- (\x+0.15,\y) -- (\x,\y-0.15) -- cycle;
            }
            
            % Деления на оси X
            \foreach \x in {0,1,2,3,4,5,6,7,8}{
                \node at (\x,-0.5) {\small \x};
            }
            
            % Деления на оси Y
            \foreach \y in {0.0,1.0}{
                \node at (-0.5,\y) {\small \y};
            }
            
            % Легенда (размещена аккуратно наверху)
            \node[gray] at (1.5,1.7) {\footnotesize Вход};
            \node[blue] at (3.8,1.7) {\footnotesize MaxPool (Сдвиг-0)};
            \node[red] at (6.8,1.7) {\footnotesize MaxPool (Сдвиг-1)};
        \end{tikzpicture}
        
        \vspace{0.5cm}
        \begin{alertblock}{Ключевое наблюдение}
            Сдвиг входного сигнала на 1 пиксель кардинально меняет результат MaxPooling операции — это и есть проявление алиасинга в CNN
        \end{alertblock}
    \end{center}
\end{frame}

% ─────────── 5–7 (НОВЫЕ) ───────────
\begin{frame}{Субдискретизация (stride $>1$)}
    \begin{block}{Определение}
        \textbf{Субдискретизация} — уменьшение пространственного разрешения за счёт выборки каждого $s$-го элемента.
    \end{block}

    \begin{columns}[T]
        \begin{column}{0.48\textwidth}
            \textbf{Где применяется}
            \begin{itemize}
                \item Max-/Average-pooling (\texttt{stride}$>1$)
                \item Свертки с \texttt{stride}$>1$
                \item Любые down-sampling-операции
            \end{itemize}
        \end{column}
        \begin{column}{0.48\textwidth}
            \textbf{Зачем нужна}
            \begin{itemize}
                \item Снижение вычислительных затрат
                \item Расширение рецептивного поля
                \item Борьба с переобучением
                \item Иерархия признаков
            \end{itemize}
        \end{column}
    \end{columns}

    \begin{alertblock}{Ключевая проблема}
        Потеря точного положения признаков → возможное нарушение инвариантности к сдвигам.
    \end{alertblock}
\end{frame}

\begin{frame}{Max-pooling \boldmath$2\times2$: визуальный пример}
    \centering
    \begin{tikzpicture}[x=0.65cm,y=0.65cm]
        % вход 4x4
        \foreach \i in {0,...,3}{
            \foreach \j in {0,...,3}{
                \draw (\i,-\j) rectangle ++(1,-1);
            }
        }
        \node at (2,-4.8) {Вход $4\times4$};

        % окно 2x2
        \draw[thick,red] (0,-0) rectangle ++(2,-2);

        % стрелка
        \draw[->,thick] (4.5,-2) -- (6.5,-2);

        % выход 2x2
        \foreach \i in {0,1}{
            \foreach \j in {0,1}{
                \draw (7+\i,-1-\j) rectangle ++(1,-1);
            }
        }
        \node at (8,-4.8) {Выход $2\times2$};
    \end{tikzpicture}

    \vspace{0.4cm}
    \begin{block}{Что происходит}
        Из каждого окна $2\times2$ берётся максимум → размер уменьшается в 4 раза, но позиционная информация стирается.
    \end{block}
\end{frame}

\begin{frame}{Сдвиг на 1 пиксель → алиасинг}
    \begin{columns}[T]
        \begin{column}{0.52\textwidth}
            \textbf{Сценарий}
            \begin{enumerate}
                \item Кот в ячейке $(10,15)$.
                \item Сдвиг всего изображения на $+1$ пиксель → кот в $(11,15)$.
                \item Применяем Max-pooling $2\times2$.
            \end{enumerate}

            \vspace{0.2cm}
            \textbf{Результат}
            \begin{itemize}
                \item Окна покрывают \emph{другие} части объекта.
                \item Активации резко меняются.
                \item Нарушается инвариантность даже при малых сдвигах.
            \end{itemize}
        \end{column}
        \begin{column}{0.46\textwidth}
            \begin{alertblock}{Последствия}
                \begin{itemize}
                    \item \textbf{Классификация:} сбой метки (кот → собака).
                    \item \textbf{Детекция:} дрейф координат, «прыгающие» боксы.
                    \item \textbf{Онлайн-системы:} нестабильные прогнозы кадр-к-кадру.
                \end{itemize}
            \end{alertblock}
        \end{column}
    \end{columns}

    \vspace{0.3cm}
    Корень проблемы — субдискретизация без антиалиасинговой фильтрации.
\end{frame}

\begin{frame}{Инвариантность и эквивариантность}
    \begin{block}{Инвариантность к сдвигам}
        Выход функции \textbf{не изменяется} при сдвиге входа
        $$f(\text{shift}(x)) = f(x)$$
        \textit{Пример: классификация — кот остается котом независимо от положения}
    \end{block}
    
    \begin{block}{Эквивариантность к сдвигам}
        Сдвиг входа приводит к \textbf{соответствующему сдвигу} выхода
        $$f(\text{shift}(x)) = \text{shift}(f(x))$$
        \textit{Пример: детекция — сдвиг объекта → сдвиг ограничивающей рамки}
    \end{block}
\end{frame}

\begin{frame}{Почему это важно и где нарушается?}
    \begin{columns}
        \begin{column}{0.5\textwidth}
            \textbf{Почему это важно?}
            \begin{itemize}
                \item Стабильность предсказаний
                \item Лучшая обобщающая способность
                \item Робастность к шуму
                \item Консистентность результатов
            \end{itemize}
        \end{column}
        \begin{column}{0.5\textwidth}
            \textbf{Где нарушается?}
            \begin{itemize}
                \item Max-pooling с stride > 1
                \item Свертки с stride > 1
                \item Average-pooling с stride > 1
                \item Любая субдискретизация
            \end{itemize}
        \end{column}
    \end{columns}
    
    \vspace{1cm}
    
    \begin{alertblock}{Ключевая проблема}
        Современные CNN архитектуры систематически нарушают инвариантность к сдвигам из-за операций субдискретизации без предварительной фильтрации
    \end{alertblock}
\end{frame}



\begin{frame}{Цель и задачи работы}
    \begin{block}{Цель}
        Разработка и анализ методов повышения инвариантности к сдвигам для CNN в задачах классификации изображений и детекции объектов
    \end{block}
    
    \textbf{Основные задачи:}
    \begin{enumerate}
        \item Теоретический анализ причин нарушения инвариантности к сдвигам
        \item Реализация и адаптация методов анти-алиасинга (BlurPool и TIPS)
        \item Разработка методики оценки инвариантности к сдвигам
        \item Сравнительный анализ методов для разных архитектур и задач
        \item Формулировка практических рекомендаций
    \end{enumerate}
\end{frame}

\section{Теоретические основы}

\section{Методы решения}

\begin{frame}{BlurPool: Антиалиасинг через фильтрацию}
    \begin{block}{Принцип работы}
        Применение низкочастотного фильтра перед субдискретизацией:
        $$\text{BlurPool}_{m,s}(x) = \text{Subsample}_{s}(\text{Blur}_{m}(x))$$
    \end{block}
    
    \begin{columns}
        \begin{column}{0.5\textwidth}
            \textbf{Используемые фильтры:}
            \begin{itemize}
                \item Triangle-3: $\frac{1}{16}[1, 2, 1]^T \cdot [1, 2, 1]$
                \item Binomial-5: $\frac{1}{256}[1, 4, 6, 4, 1]^T \cdot [1, 4, 6, 4, 1]$
            \end{itemize}
            
            \textbf{Модификации операций:}
            \begin{itemize}
                \item $\text{MaxPool}_{k,s} \rightarrow \text{Subsample}_{s} \circ \text{Blur}_{m} \circ \text{Max}_{k,1}$
                \item $\text{Conv}_{k,s} \rightarrow \text{Subsample}_{s} \circ \text{Blur}_{m} \circ \text{Conv}_{k,1}$
            \end{itemize}
        \end{column}
        \begin{column}{0.5\textwidth}
            \textbf{Преимущества:}
            \begin{itemize}
                \item Простота интеграции
                \item Минимальные вычислительные затраты (<1\%)
                \item Применимо к предобученным моделям
                \item Значительное улучшение инвариантности
            \end{itemize}
        \end{column}
    \end{columns}
\end{frame}

\begin{frame}{TIPS: Полифазная выборка}
    \begin{block}{Принцип работы}
        Разбиение сигнала на $s^2$ фаз с последующим взвешенным объединением:
        $$\text{TIPS}_s(x) = \sum_{i=0}^{s-1}\sum_{j=0}^{s-1} \tau_{is+j} \cdot \text{Subsample}_{s}(\text{Shift}_{(i,j)}(x))$$
    \end{block}
    
    \begin{columns}
        \begin{column}{0.5\textwidth}
            \textbf{Особенности:}
            \begin{itemize}
                \item $s^2$ параллельных вычислительных путей
                \item Обучаемые смешивающие коэффициенты $\tau_{is+j}$
                \item Теоретическая гарантия инвариантности
                \item Обработка субпиксельных сдвигов
            \end{itemize}
        \end{column}
        \begin{column}{0.5\textwidth}
            \textbf{Преимущества и недостатки:}
            \begin{itemize}
                \item[+] Максимальная инвариантность
                \item[+] Теоретические гарантии
                \item[+] Высокая стабильность
                \item[-] Большие вычислительные затраты
                \item[-] Сложность интеграции
            \end{itemize}
        \end{column}
    \end{columns}
\end{frame}

\section{Экспериментальная часть}

\begin{frame}{Экспериментальная методология}
    \begin{block}{Используемые датасеты}
        \begin{itemize}
            \item \textbf{CIFAR-10} — классификация (60,000 изображений 32×32)
            \item \textbf{ImageNet-mini} — классификация (100 классов, 1300 изображений)
            \item \textbf{Imagenette} — быстрые эксперименты (10 классов)
            \item \textbf{COCO-sample} — детекция объектов
        \end{itemize}
    \end{block}
    
    \begin{block}{Тестирование инвариантности}
        \begin{itemize}
            \item Синтетические последовательности сдвигов (0-8 пикселей)
            \item Субпиксельные сдвиги через билинейную интерполяцию
            \item Комбинированные сцены для детекции
            \item До 100 прогонов для каждого изображения
        \end{itemize}
    \end{block}
\end{frame}

\begin{frame}{Архитектуры и модификации}
    \begin{columns}
        \begin{column}{0.5\textwidth}
            \textbf{Классификационные модели:}
            \begin{itemize}
                \item VGG16 — замена max-pooling слоев
                \item ResNet50 — модификация сверток с stride=2
                \item Применимо к предобученным моделям
            \end{itemize}
            
            \vspace{0.5cm}
            
            \textbf{Детектор объектов:}
            \begin{itemize}
                \item YOLOv5s — сложная архитектура
                \item Backbone (CSPDarknet)
                \item Neck (PANet)
                \item Head — предсказания
            \end{itemize}
        \end{column}
        \begin{column}{0.5\textwidth}
            \textbf{Метрики оценки:}
            
            \textit{Классификация:}
            \begin{itemize}
                \item Top-1 Accuracy
                \item Stability (при сдвигах 0-5 пикс.)
                \item Consistency
                \item Flip-Invariance
            \end{itemize}
            
            \textit{Детекция:}
            \begin{itemize}
                \item mAP (mean Average Precision)
                \item IoU Stability
                \item Center Drift (дрейф центра)
                \item Size Consistency
            \end{itemize}
        \end{column}
    \end{columns}
\end{frame}

\section{Результаты}

\begin{frame}{Результаты для классификации}
    \begin{table}[h]
    \centering
    \small
    \begin{tabular}{|l|c|c|c|c|}
    \hline
    \textbf{Модель} & \textbf{Top-1 Acc.} & \textbf{Stability} & \textbf{Consistency} & \textbf{Flip-Inv.} \\
    \hline
    VGG16 & 0.7135 & 0.8612 & 0.9103 & 0.8821 \\
         VGG16 + BlurPool & 0.7262 & \textcolor{darkgreen}{0.9358} & \textcolor{darkgreen}{0.9642} & 0.9127 \\
     VGG16 + TIPS & 0.7298 & \textcolor{darkgreen}{0.9687} & \textcolor{darkgreen}{0.9801} & 0.9236 \\
    \hline
    ResNet50 & 0.7613 & 0.8874 & 0.9256 & 0.9014 \\
         ResNet50 + BlurPool & 0.7725 & \textcolor{darkgreen}{0.9432} & \textcolor{darkgreen}{0.9715} & 0.9325 \\
     ResNet50 + TIPS & 0.7769 & \textcolor{darkgreen}{0.9742} & \textcolor{darkgreen}{0.9857} & 0.9412 \\
    \hline
    \end{tabular}
    \end{table}
    
    \begin{block}{Ключевые результаты}
        \begin{itemize}
            \item Улучшение стабильности на \textbf{7.5-10.7 п.п.}
            \item TIPS показывает лучшие результаты по всем метрикам
            \item Одновременное улучшение точности и стабильности
        \end{itemize}
    \end{block}
\end{frame}

\begin{frame}{Результаты для детекции объектов}
    \begin{table}[h]
    \centering
    \small
    \begin{tabular}{|l|c|c|c|c|}
    \hline
    \textbf{Модель} & \textbf{mAP} & \textbf{IoU Stability} & \textbf{Center Drift} & \textbf{Size Cons.} \\
    \hline
    YOLOv5s & 0.371 & 0.853 & 4.87 & 0.912 \\
         YOLOv5s + BlurPool & 0.384 & \textcolor{darkgreen}{0.921} & \textcolor{darkgreen}{2.31} & 0.958 \\
     YOLOv5s + TIPS & 0.392 & \textcolor{darkgreen}{0.964} & \textcolor{darkgreen}{1.45} & 0.974 \\
    \hline
    \end{tabular}
    \end{table}
    
    \vspace{0.3cm}
    
    \begin{columns}
        \begin{column}{0.5\textwidth}
            \textbf{Основные улучшения:}
            \begin{itemize}
                \item IoU Stability: +6.8-11.1 п.п.
                \item Center Drift: снижение в 2+ раза
                \item mAP: улучшение на 1.3-2.1 п.п.
            \end{itemize}
        \end{column}
        \begin{column}{0.5\textwidth}
            \textbf{Особенно для малых объектов:}
            \begin{itemize}
                \item Наибольший прирост точности
                \item До 2.6 п.п. для TIPS
                \item Критично для практических задач
            \end{itemize}
        \end{column}
    \end{columns}
\end{frame}

\begin{frame}{Вычислительная эффективность}
    \begin{columns}
        \begin{column}{0.5\textwidth}
                        \textbf{Классификационные модели:}
            \begin{table}[h]
            \centering
            \footnotesize
            \begin{tabular}{|l|c|c|}
            \hline
            \textbf{Модель} & \textbf{FLOPs (G)} & \textbf{Overhead} \\
            \hline
            VGG16 & 15.47 & - \\
             + BlurPool & 15.58 & \textcolor{darkgreen}{0.7\%} \\
             + TIPS & 15.83 & \textcolor{orange}{2.3\%} \\
            \hline
            ResNet50 & 4.12 & - \\
             + BlurPool & 4.17 & \textcolor{darkgreen}{1.2\%} \\
             + TIPS & 4.32 & \textcolor{orange}{4.9\%} \\
            \hline
            \end{tabular}
            \end{table}
        \end{column}
        \begin{column}{0.5\textwidth}
                        \textbf{Детекция объектов:}
            \begin{table}[h]
            \centering
            \footnotesize
            \begin{tabular}{|l|c|c|}
            \hline
            \textbf{Модель} & \textbf{FPS} & \textbf{Overhead} \\
            \hline
            YOLOv5s & 47.3 & - \\
             + BlurPool & 45.8 & \textcolor{darkgreen}{1.2\%} \\
             + TIPS & 45.1 & \textcolor{orange}{3.6\%} \\
            \hline
            \end{tabular}
            \end{table}
        \end{column}
    \end{columns}
    
    \vspace{0.5cm}
    
    \begin{block}{Вывод}
        \begin{itemize}
            \item BlurPool: минимальные затраты (<1\% FLOPs)
            \item TIPS: умеренные затраты (2-5\% FLOPs)
            \item Пригодны для систем реального времени
        \end{itemize}
    \end{block}
\end{frame}

\section{Практические рекомендации}

\begin{frame}{Рекомендации по применению}
    \begin{block}{Системы критического применения}
        \textbf{Рекомендация:} TIPS
        \begin{itemize}
            \item Медицинская диагностика, автономное вождение
            \item Максимальная стабильность (97-98\% консистентности)
            \item Минимальный дрейф локализации
            \item Оправданы дополнительные 4-5\% вычислений
        \end{itemize}
    \end{block}
    
    \begin{block}{Системы с ограниченными ресурсами}
        \textbf{Рекомендация:} BlurPool
        \begin{itemize}
            \item Мобильные приложения, встраиваемые системы
            \item Существенное улучшение (91-94\% консистентности)
            \item Минимальные затраты (<1\% FLOPs)
            \item Идеально для батарейных устройств
        \end{itemize}
    \end{block}
    
    \begin{block}{Детекция малых объектов}
        Особенно рекомендуется применение методов антиалиасинга
        \begin{itemize}
            \item Наибольший выигрыш в точности (до 2.6 п.п.)
            \item Критично для систем видеонаблюдения
        \end{itemize}
    \end{block}
\end{frame}

\section{Заключение}

\begin{frame}{Основные результаты}
    \begin{block}{Научная новизна}
        \begin{itemize}
            \item Комплексный анализ инвариантности к сдвигам в CNN
            \item Первое детальное исследование влияния на детекторы объектов
            \item Математическая формализация проблемы алиасинга в CNN
            \item Адаптация методов антиалиасинга для YOLOv5
        \end{itemize}
    \end{block}
    
    \begin{block}{Практическая значимость}
        \begin{itemize}
            \item Повышение стабильности CNN на 40-90\%
            \item Улучшение точности классификации на 1.3-2.1 п.п.
            \item Снижение дрейфа локализации в 2+ раза
            \item Минимальные вычислительные затраты
            \item Простота интеграции в существующие системы
        \end{itemize}
    \end{block}
\end{frame}

\begin{frame}{Выводы}
    \begin{enumerate}
        \item \textbf{Проблема подтверждена:} Современные CNN демонстрируют значительную нестабильность при субпиксельных сдвигах (до 30\% вариации для классификации, 25\% для детекции)
        
        \item \textbf{Методы эффективны:} BlurPool и TIPS обеспечивают существенное повышение инвариантности при минимальных вычислительных затратах
        
        \item \textbf{Практическая применимость:} Методы легко интегрируются в существующие архитектуры без значительного переобучения
        
        \item \textbf{Рекомендации сформулированы:} Даны четкие критерии выбора методов для различных типов приложений
        
        \item \textbf{Перспективы развития:} Направления для дальнейших исследований в области робастности CNN
    \end{enumerate}
\end{frame}

\begin{frame}
    \begin{center}
        \Huge Спасибо за внимание!
    \end{center}
\end{frame}

\section{Дополнительные слайды}

\begin{frame}{Детальный анализ стабильности}
    \begin{columns}
        \begin{column}{0.5\textwidth}
            \textbf{Метрики стабильности:}
            \begin{itemize}
                \item Consistency: $\frac{1}{N}\sum_{i=1}^N \cos(f(x_i), f(x_i'))$
                \item IoU Stability: средняя IoU при сдвигах
                \item Center Drift: $\sqrt{(\Delta x)^2 + (\Delta y)^2}$
            \end{itemize}
        \end{column}
        \begin{column}{0.5\textwidth}
            \textbf{Экспериментальный протокол:}
            \begin{itemize}
                \item 8 направлений сдвига
                \item Субпиксельные сдвиги 0.1-0.9
                \item Статистическая значимость p<0.01
                \item Воспроизводимые результаты
            \end{itemize}
        \end{column}
    \end{columns}
\end{frame}

\begin{frame}{Сравнение с другими методами}
    \begin{table}[h]
    \centering
    \small
    \begin{tabular}{|l|c|c|c|}
    \hline
    \textbf{Метод} & \textbf{Стабильность} & \textbf{Вычисл. затраты} & \textbf{Интеграция} \\
    \hline
    Стандартные CNN & Низкая & Базовая & - \\
    Data Augmentation & Умеренная & Высокие & Легкая \\
    BlurPool & Высокая & Минимальные & Легкая \\
    TIPS & Максимальная & Умеренные & Средняя \\
    \hline
    \end{tabular}
    \end{table}
    
    \begin{block}{Преимущества предложенных методов}
        \begin{itemize}
            \item Архитектурный подход vs. аугментация данных
            \item Гарантированная инвариантность vs. случайная стабилизация
            \item Эффективность на этапе инференса
        \end{itemize}
    \end{block}
\end{frame}

\end{document}